\chapter*{Conclusion and future work}
\addcontentsline{toc}{chapter}{Conclusion and future work}

\section*{Conclusion}

In this thesis, we developed a Mathlib-compatible framework for formal
group laws in Lean.  Our main results can be summarized as follows.

First, we formalized formal group laws over a commutative ring together
with their commutativity, homomorphisms, and isomorphisms, providing a
foundational API for working with formal group laws in Lean.

Second, we proved that a formal group law induces a group structure on
its points, which is commutative whenever the law is.  This shows that
the algebraic structure expected of formal group laws can be recovered
and used formally.

Third, we formalized the functional equation lemma, including both the
construction satisfying the functional equation and the descent
statement that the resulting formal group law has coefficients in the
required subring.

Along the way, we established a number of supporting results about
multivariate power series, which provide reusable infrastructure for
future work involving power series in Mathlib.  Finally, we formalized
the height of a formal group law, bridging the basic algebraic theory to
deeper classification results and arithmetic applications.

\section*{Future Work}

The framework developed in this thesis opens up several natural
directions for further work.

The most immediate is Lubin--Tate theory, which uses formal group laws to
construct highly structured abelian extensions of local fields.
Formalizing it would combine our formal group law API with Mathlib's
existing theory of local fields, valuations, and ramification, and would
be a significant step toward explicit local class field theory.

A second direction concerns the formal group laws of elliptic curves.
Every elliptic curve has an associated formal group law, obtained by
completing the curve at the identity.  Over a field of positive
characteristic, the height of this law records important arithmetic
information: it equals $1$ in the ordinary case and $2$ in the
supersingular case.  Formalizing this link would connect our theory to
Mathlib's existing formalization of elliptic curves.

A third goal is Lazard's theorem, which asserts the existence of a
universal formal group law and identifies its coefficient ring as a
polynomial ring.  Fundamental in the algebraic theory of formal group
laws and central to complex cobordism in algebraic topology, its
formalization would require a careful development of the Lazard ring and
the accompanying universal constructions.

A fourth direction is the classification of formal group laws by height,
the central invariant in positive characteristic.  A natural long-term
goal is to formalize the classification of formal group laws of a given
height, including the construction of isomorphisms after a suitable base
change, thereby extending the basic theory of this thesis to a more
structural understanding.

Taken together, these directions build on the reusable infrastructure
developed in this thesis as a foundation for further work in arithmetic
geometry.